\section{Related Works}
\subsection{LLM Agents}
Recent work has extended the capabilities of LLMs beyond standalone reasoning and understanding, enabling them to operate as multi-agents that can interact with environments, tools, and other agents to perform complex tasks \citep{wu2023autogen, camel, agentverse, chatdev, xu2025instructagent, xu2025mem, he2025you, le2025instructiontuningcotprompting, 11079788, sym17071087,yu2025gradient,cheng2024optimized,wang2024made, ma2024comparative, song2024u2++}. These multi-agent systems (MAS) integrate various techniques, including CoT prompting \citep{wei2022chain,zhang2024cut}, iterative refinement \citep{shinn2023reflexion}, self-improvement \citep{huang2023large,mei2025omnirouter}, and external tool usage \citep{hao2023toolkengpt, qintoolllm, schick2023toolformer,zhang2025multi,Shi2024,zhu2025llm}, to support multi-step decision-making and long-horizon planning. They have been applied successfully in domains such as mathematical reasoning \citep{agentverse}, software engineering \citep{chatdev, jimenez2023swe, yang2024swe, wang2024opendevin}, and scientific discovery \citep{ai_scientist, agent_laboratory}. Agent frameworks typically structure the interaction with LLMs using techniques such as few-shot prompting \citep{brown2020language, sun2025cross} and guided reasoning \citep{wei2022chain,shinn2023reflexion,Lorasculpt_CVPR25,huang2025keeping,sun2023prompt}, relying on the model’s in-context learning capabilities \citep{olsson2022context,jin2025massive}.

\subsection{Test-time Scaling}
A wide range of methods have been developed to improve reasoning in LLMs by leveraging test-time scaling (TTS). Recent work explores techniques including hierarchical hypothesis search, which enables inductive reasoning through structured exploration \citep{wang2023hypothesis}, and tool augmentation during inference, which enhances downstream performance by allowing models to interact with external environments \citep{gao2023pal, qintoolllm,guo2025reagan}. Other approaches focus on internal mechanisms, such as learning thought tokens in an unsupervised manner \citep{zelikman2024quiet, goyal2023think}, allowing models to better utilize extended reasoning sequences. Among the most studied scaling paradigms are parallel and sequential TTS approaches. Parallel methods generate multiple solution candidates independently and select the best one using a scoring criterion, such as majority voting or outcome-based reward models \citep{brown2024large, irvine2023rewarding, snell2025scaling}. In contrast, sequential methods condition each new attempt on the previous ones, allowing iterative refinement based on prior outputs \citep{s1,snell2025scaling, hou2025advancing, lee2025evolving, cai2024role, shi2024scaling, he2025self,zhang2025growing}. Bridging these strategies, tree-based techniques such as Monte Carlo Tree Search (MCTS) \citep{zhang2023planning, zhou2023language, zhang2024llama, wu2024inference} and guided beam search \citep{xie2023self} enable structured exploration through branching and evaluation. Central to many of these methods are reward models, which provide feedback signals for generation. These can be categorized as outcome reward models, which evaluate entire solutions \citep{xin2024deepseek, ankner2024critique}, or process reward models, which assess intermediate reasoning steps \citep{lightman2023let, wang2023math, wu2024inference}, guiding the model toward more effective reasoning paths.